\cotefile{156-7}
\scanpage{34}
\mspage{33}

\uline{Lemme} 8 \deleted{\unsure{Soit}}\unsure{Soient} $X,Y\in M$. Alors

\margin{$\mathrm{At\,L\,2}$}
\[
X\,\underbrace{\mathrel{\unsure{\ooalign{\hfil$>$\hfil\cr\hfil$<$\hfil\cr}|\gtrless}}Y
  \text{ et }\widetilde X\cap\widetilde Y}
 _{\substack{\text{i.e. }X\parallel Y\\\text{au sens de l'}\\\text{atlas}}}=\varnothing
\quad\Longleftrightarrow\quad
\begin{aligned}
 &\forall\,x\in\unsure{\mathrm{omb}|\mathrm{ends}|\mathrm{ens}}(\unsure{\alpha|X})\\[-2pt]
 &\phantom{\forall\,}y\in\mathrm{omb}(Y)
\end{aligned}
\quad\Longrightarrow\quad x\parallel y
\]

C'est tautologique !

\deleted{\uline{\unsure{Lemme}} \unsure{9}} Ainsi, \unsure{notre} atlas satisfait, en plus \unsure{de|des} $\mathrm{At\,1}-\mathrm{At\,7}$, \deleted{\unsure{la}}les conditions sur les lieux $\mathrm{At\,L\,1}$, $\mathrm{At\,L\,2}$, $\mathrm{At\,L\,3}$. Reste à voir la condition $\mathrm{At\,CL(spac)}$ -- mais elle est vraie par définition.

On trouve ainsi le

\uline{Théorème-\unsure{Scholie|Solution}} Soit $(M,{\leq},{\ll})$ un ensemble muni \unsure{de|des} deux relations d'ordre, $\unsure{L|\mathcal L}$ l'ens des éléments minimaux de $M$ pour $\ll$, $\parallel$ une relation symétrique antiréflexive sur $L$. Pour que cet \inserted{ens de} données \deleted{\unsure{provienne}} soit isom. à celui provenant d'un atlas \inserted{spacieux} satisfaisant \deleted{$\mathrm{At\,1}$} (en plus de $\mathrm{At\,1}-\mathrm{At\,7}$) les axiomes sur les lieux $\mathrm{At\,L\,1}$, $\mathrm{At\,L\,2}$, $\mathrm{At\,L\,3}$, \deleted{\unsure{p}}plus $\mathrm{At\,\deleted{\unsure{L}}CL(spac)}$, il faut \unsure{s.} que les conditions suivantes soient satisfaites :
\scanpage{35}
\mspage{34}

% slanted
\margin{axiome trivial}
\begingroup\makeatletter\tagsleft@true\makeatother
\[
 X\leq Y\quad\Longrightarrow\quad X\ll Y
 \tag*{$(\deleted{\unsure{C|G}}\mathrm{ML\,0})$}
\]

% slanted
\margin{axiome du raffinement induit}
$(\deleted{C}\mathrm{ML\,1})$ \deleted{\unsure{ent}}\unsure{Soit|Soient} $\deleted{\unsure{X'}}$
\diagram{\begin{tikzcd}[ampersand replacement=\&,row sep=0pt,column sep=3pt]
 Z \& \ll \& X\\
 \& \& \rotatebox{90}{$\leq$}\\
 \& \& X'
\end{tikzcd}}
dans $M$, et soit $x\in\unsure{L|\mathcal L}$ avec $x\ll Z$, $x\ll X'$. Alors $\exists\,Z'\in M$ avec
\diagram{\begin{tikzcd}[ampersand replacement=\&,row sep=1pt,column sep=4pt]
 \& \& Z \& \ll \& \unsure{x|X}\\
 \& \& \rotatebox{90}{$\leq$} \& \& \rotatebox{90}{$\leq$}\\
 x \& \ll \& Z' \& \ll \& X'\rlap{$\quad\unsure{,|.}$}
\end{tikzcd}}

\uline{De} plus, si $Z\leq X$, on peut supposer $Z'\leq X'$.

% slanted
\margin{axiome de divisibilité \deleted{\unsure{pour|par}}}
\[
\begin{aligned}
 &\forall X\in M,\quad \unsure{\mathrm{omb}|\mathrm{ends}|\mathrm{ens}}(X)^{\unsure{\circ|0}}\ne 0
   \quad\text{i.e. }\exists\,x\in L,\ \text{avec}\\
 &x\mathrel{\overset{\circ}{\ll}}X\quad\text{i.e. }x\ll X\text{ et }
   (\unsure{a|x|\alpha}\ll X'\leq X_{\unsure{x|n}})\Longrightarrow X'=X
\end{aligned}
\tag*{$(\mathrm{ML\,2})^{*}$}
\]

% slanted
\margin{axiome de \deleted{\unsure{plus}} la \deleted{plus \unsure{petite}} strate minimale}
\[
\begin{aligned}
 &\deleted{\unsure{\forall x\in}}\ \forall\,X\in M,\ x\unsure{\in|\notin}L
    \ \text{avec }\ x\ll X,\ \exists\,X'\leq X\\
 &\text{avec }\quad x\mathrel{\overset{\circ}{\ll}}X'\,.
\end{aligned}
\tag*{$(\mathrm{ML\,3})$}
\]

$(\mathrm{ML\,4})$ $\parallel$ est \unsure{sym,|sym.} antiréflexive \inserted{$(x\parallel y\Rightarrow y\parallel x\text{ et }x\ne y)$}

% slanted
\margin{C'est le seul des axiomes qui \unsure{fasse} intervenir la relation $\parallel$}
\margin{$\swarrow\!\big\lfloor$ axiome de disjonction pour les lieux}
\[
\left\{\begin{aligned}
 &(\mathrm{ML}\,\deleted{\unsure{4}}5)\qquad \forall\,X,Y\in M\ \text{avec }\ Y\leq X,\quad\deleted{\text{\unsure{si }}y}\\
 &\hspace{3.5em} y\in\mathrm{omb}(Y),\ x\in\unsure{\mathrm{om}|\mathrm{en}|\mathrm{omb}}(X)
    \setminus\mathrm{omb}(Y)\quad\Longrightarrow\quad x\parallel y
\end{aligned}\right.
\]
\endgroup

\deleted{Les atlas correspondant à une telle donnée, réalisés comme des ens de parties de $M$}
\[
\deleted{\unsure{\mathcal F|F}\subset\unsure{\mathrm P|\mathcal P}(M),}
\]
\deleted{sont \deleted{\unsure{toute|toutes} \unsure{forme}} les $\mathcal F$ satisfaisant les conditions suivantes}

\deleted{(i)}

$*$ \uline{NB} Cet axiome implique que si \deleted{\unsure{la}}les relations d'ordre $\leq$ et $\ll$ sont égales, alors $\deleted{\unsure{X}}M$ est discret i.e. $M=L$. \deleted{\unsure{La thèse est vide}} L'axiome $\mathrm{ML\,4}$ est vide !
\scanpage{36}
\mspage{35}

Soit $\deleted{\unsure{\mathcal F_{\mathrm{max.}}}\ \unsure{\mathcal F^{*}}}$ $\Phi$ l'ens des parties de $M$ qui sont a) fermées pour $\leq$ et b) formées d'éléments deux à deux compatibles. (au sens de (1) p.~29). Munissons $\Phi$ de la relation $\leq\ \overset{\mathrm{def}}{=}$ l'inclusion, et $\ll$ définie par $F\ll G\ \overset{\mathrm{def}}{\Longleftrightarrow}\ \forall X\in F,\exists\,Y\in G$ avec $X\ll Y$. Alors $\Phi$ est un atlas \deleted{admettant} satisfaisant $\mathrm{\deleted{\unsure{M|A}}At\,L\,1}$, $\mathrm{At\,L\,2}$, $\mathrm{At\,L\,3}$ et $\mathrm{At\,CL(\unsure{\mathrm{spac}|\mathrm{spa}})}$, et admettant les $\widetilde X=M_{\leq\unsure{X|x}}$ $\deleted{\unsure{\{\{X\}\}|\{\{x\}\}}}$ $(X\in M)$ comme \deleted{objets} figures élémentaires, \deleted{les $\{\{x\}\}$ \unsure{pour}} et \deleted{\unsure{qui}} induisant sur $M$ les relations $\leq$, $\ll$ de départ, et sur $\unsure{L|\mathcal L}$ \inserted{(pour la \deleted{\unsure{dis}} disjonction des figures)} la relation $\parallel$.

\deleted{Tout \unsure{autre} \unsure{les atlas|tel atlas}} \inserted{Tout atlas correspondant à ces} $\unsure{\Rightarrow|\mathord{=}}$ données s'identifie à \unsure{un|son} \uline{sous-atlas} \uline{spacieux} de $\Phi$\unsure{, »|,} \inserted{\unsure{i.e.|i.e}} \deleted{(parties $\unsure{\mathcal F|F}$ de $\Phi$} à une partie
\[
 \mathcal F\subset\Phi
\]
satisfaisant

(\deleted{\unsure{a|+}}i) $\forall\,F\in\mathcal F$, et $G\leq F$ \deleted{\unsure{sous-figure}} \inserted{(partie fermée} de $F,\leq$) \unsure{on a} $G\in\mathcal F$ et

(ii) $\forall\,F,G\in\mathcal F$, Si $F\mathrel{\unsure{\ooalign{\hfil$>$\hfil\cr\hfil$<$\hfil\cr}|\gtrless}}G$ i.e. $F\cup G\in\Phi$ (i.e. \deleted{\unsure{tous}}\unsure{tout|tous} éléments de $F_{\unsure{a|x}}$ est comp. à tout \unsure{élément|élmt|élémts} de $G$) $F\cup G\in\mathcal F$.
\scanpage{37}
\mspage{36}

Il contient de plus les \deleted{objets} figures élémentaires, i.e. satisfait :
\begingroup\makeatletter\tagsleft@true\makeatother
\[
 \forall\,\unsure{x|X}\in M,\quad \{\{X\}\}\in\unsure{\mathcal F|F}
 \tag*{(iii)}
\]
\endgroup
et par suite, il contient toutes les parties fermées de réunions finies de tels $\widetilde X$.

Inversement, \unsure{tte} $\mathcal F$ satisfaisant ces trois conditions satisfait aux conditions voulues

\deleted{Si \unsure{les}} \inserted{Pour que les} figures de $\mathcal F$ soient de \unsure{t.f.}, \unsure{il f} et s. qu'on \deleted{ait}
\begingroup\makeatletter\tagsleft@true\makeatother
\[
 \mathcal F=\mathcal F_{0}=
 \left\{\begin{aligned}
 &\text{réunions finies d'ensembles}\\
 &\qquad\text{de la forme }\widetilde X
 \end{aligned}\right\}
 \tag*{a)}
\]
\endgroup

\unsure{(ce|Ce} qui ne donne un atlas que si on suppose de plus

b) toute partie fermée d'un $\widetilde X$ \inserted{(pour $\leq$)} est de type fini

À part donc un petit flottement dû à la possibilité de \deleted{\unsure{feuilles|feuillets}} \inserted{figures} \deleted{\unsure{finies}} infinies, \deleted{\unsure{o}}\unsure{on} voit que \unsure{tout|\deleted{tout}} s'exprime par $(M,{\leq},{\ll},{\parallel})$, avec $\parallel$ \unsure{définie|défini} sur $\unsure{L|\mathcal L}$ sans plus, \deleted{\unsure{(et si on dit aussi|(et sinon de aussi|(et sinon de deux}} satisfaisant aux quatre axiomes (cinq avec $(\mathrm{ML\,0})$) $\mathrm{ML\,1}-\mathrm{ML\,4}$ (au lieu
\scanpage{38}
\mspage{37}

de la douzaine d'axiomes $\deleted{\unsure{p}}\mathrm{At\,1}-7$, $\mathrm{At\,L\,1}-3$, $\mathrm{At\,\deleted{\unsure{L}}CL(spac)}$.

Si on se donne $\parallel$ sur $M$, et non sur $\unsure{L|\mathcal L}$, il faut simplement rajouter un axiome supplémentaire :

$(\mathrm{ML}\,\deleted{\unsure{5}}6)$ Soit $\unsure{X|x},Y\in\deleted{\unsure{L}}M$. Alors
\[
 X\parallel Y\quad\Longleftrightarrow\quad \forall\,x\unsure{\in|,}y\deleted{\unsure{-}}\in L,\quad
 x\ll X,\ y\deleted{\unsure{\leq}}\ll Y\ \Rightarrow\ x\parallel y.
\]

\uline{Variante pour atlas polyédraux}. \inserted{(satisfaisant $\mathrm{At\,1}-\mathrm{At\,7}$, $\mathrm{At\,L\,1}-\mathrm{At\,L\,3}$, \inserted{$\mathrm{At\,CL(pol)}$})} C'est exactement la \unsure{$=$} chose, sauf qu'on change en conséquence la définition de $X\mathrel{\unsure{\ooalign{\hfil$>$\hfil\cr\hfil$<$\hfil\cr}|\gtrless}}Y$ (p.~29), et qu'il faut supposer sur $M,\leq$ une condition supplémentaire

$(\mathrm{ML\,0\ pol})$ $\deleted{\unsure{\exists}}\forall\,X\in M$, $\widetilde X=M_{\leq X}$ est un

\noindent\begin{tabular}{@{}|p{\dimexpr\linewidth-2\tabcolsep-1pt\relax}@{}}
ens ordonné \uline{polyédral} (p.~\unsure{13|15}), \unsure{i.e|ie} $\deleted{\exists\,\unsure{X|x},Y,Z}$ deux \unsure{éléments|élmts} $Y,Z$ de $\deleted{\unsure{X}}M$, \inserted{$\leq$} qui sont majorés (par un $X$) et minorés\unsure{,} \deleted{\unsure{adme}} (par un $T$) admettent un inf -- ou \unsure{encore}, les ensembles $M_{T\leq\unsure{\cdot|\bullet}\leq\unsure{x|X}}$ ont des inf finis (donc aussi des \unsure{sup|type} finis s'ils sont finis).
\end{tabular}
\scanpage{39}
\mspage{38}

Ça montre, de plus, que si on part de $(M,{\leq},{\ll})$, il y a un très large choix pour les $\parallel$ possibles sur $\unsure{L|\mathcal L}$, i.e. les parties \unsure{envisagées} de $\unsure{\mathrm P|\mathcal P}_2(L)$. Si on désigne par
\[
 \unsure{\text{``}|\mathinner{..}}\quad
 \Gamma_0=\left\{\{x,y\}\in\mathrm P_2(L)\ \middle|\
 \begin{aligned}
 &\exists\,X,Y\in M,\ Y\leq X\\
 &y\deleted{\unsure{\leq}}\ll\deleted{\unsure{X}}\unsure{Y|X},\
 x\deleted{\unsure{\leq}}\ll X,\ x\not\ll Y
 \end{aligned}\right\}
\]
on peut prendre comme structure de graphe $\Gamma$ \deleted{\unsure{de}} (pour une relation $\parallel$) n'importe quelle partie $\Gamma\supset\Gamma_0$.
\[
 \Gamma_0\subset\Gamma\subset\mathrm P_2(L)\,.
\]

Ça fait de la marge ! Pour la suite, il faudra voir comment les autres axiomes affectent le choix de $\parallel$ (et du reste). Pour le moment, je songe surtout à
\[
 \begin{aligned}
 &\mathrm{At\ supp} &&\text{\unsure{(cf|cf} page 18)}\\
 &\mathrm{At\ \unsure{ens|ends}}&&\text{et ses variantes}\quad
 \text{(cf GF$\overline{\underline{\unsure{\mathrm{VI}|\mathrm{VII}}}}$ p 52 \unsure{....|...})}
 \end{aligned}
\]

\noindent\begin{tabular}{@{}||p{\dimexpr\linewidth-2\tabcolsep-3pt\relax}@{}}
\deleted{\unsure{Dans|Pour}} Pour qu'on ait un atlas ensembliste \inserted{spacieux}, \inserted{(resp. \unsure{polyédrisé|polyédralisé|polyédral})} suffit-il que $x\parallel y\Longleftrightarrow x\ne y$ ?? Sans
\end{tabular}
doute pas, car cela n'implique aucune restriction sur les autres données $\leq$, $\ll$, et les conditions $\mathrm{ML\,0}$ à $\mathrm{ML\,3}$ n'impliquent \uline{pas} que l'application \unsure{enlocus|unilocus|endroits|endroit|ombres} : $M\longrightarrow\mathrm{Fig}(L)$ soit injective, ni que $\ll$ soit induite par la relation $\ll$ entre fig. ensemblistes \inserted{dans $L$} \unsure{$\to$}
\scanpage{40}
\mspage{39}

Dans les raisonnements précédents, le fait que les éléments de $\unsure{L|\mathcal L}$ soient minimaux n'est intervenu qu'à la fin, pour prouver que la relation de \unsure{disjoint|disjoint.} $\parallel$ sur $L$ est \unsure{bien} celle induite par la disjonction des lieux, au sens de l'atlas construit. On va tout reprendre.

\noindent\begin{tabular}{@{}|p{\dimexpr\linewidth-2\tabcolsep-1pt\relax}@{}}
\uline{Définition} \deleted{\unsure{L|S}}Soit $\Lambda\subset M$. On dit que $\Lambda$ est \uline{dense} \inserted{dans $M$} si $\forall\,X,$ \deleted{\unsure{avec}} $\in M,\ \exists\,Y\in\Lambda$ avec $Y\mathrel{\overset{\circ}{\ll}}X$.
\end{tabular}

\uline{Ex} $1^{\circ})\ \Lambda=M$, exemple ``universel'' (on prend $Y=X$)

$2^{\circ})\ \Lambda=L$ \deleted{\unsure{que}} quand $\mathrm{At\,L\,1}$ est satisfait

On va énoncer des axiomes sur $(M,{\leq},{\ll},{\unsure{\ooalign{\hfil$>$\hfil\cr\hfil$<$\hfil\cr}|\gtrless}},\Lambda\subset M)$ qui assurent que $(M,{\leq},{\ll},{\ooalign{\hfil$>$\hfil\cr\hfil$<$\hfil\cr}})$ \unsure{provienne|provient} d'un atlas, et que $\Lambda\subset M$ est dense, avec \deleted{\unsure{un|en}} en plus sur l'atlas (au besoin) la condition $\mathrm{At\,D}$ ou $\mathrm{At\,D}_{\unsure{0|o}}$ ou $\mathrm{At\,C}$.

% slanted
\margin{Axiome trivial}
\begingroup\makeatletter\tagsleft@true\makeatother
\[
 X\leq Y\ \Longrightarrow\ X\ll Y\qquad\text{comme avant}
 \tag*{$\mathrm M\Lambda\,0$}
\]

% slanted
\margin{Axiome du raffinement induit}
% margin note written sideways
\margin{$\left.\rule{0pt}{3em}\right\}$ \deleted{\unsure{peu|texte} \illegible{1} \unsure{pas|pour} \unsure{de 1}}}
$\mathrm M\Lambda\,1$
\diagram{\begin{tikzcd}[ampersand replacement=\&,row sep=0pt,column sep=3pt]
 \text{Si }\ Y\deleted{\unsure{Z|X}}\ll \& X
     \& \text{et }\deleted{\unsure{x|Y}}\ \deleted{\unsure{x|Y}}Z\mathrel{\text{\uline{$\in\Lambda$}}}\quad\text{avec}\\
 \& \rotatebox{90}{$\leq$} \&\\
 Z\ll Y,\ Z\ll X',\ \deleted{\unsure{\text{Alors},\exists}} \& X'
     \& ,\ \text{alors }\exists\,Y'\text{ avec}
\end{tikzcd}}
\[
 \begin{aligned}
 &\deleted{Z'\unsure{\leq|<}Z,}\quad
 \deleted{Z'}\unsure{Y'|Z'}\leq\deleted{Z}\unsure{Y|Z},\
 \deleted{Z'}\unsure{Y'|Z'}\ll X',\ Z\ll Y'\\
 &\text{et si }\deleted{Z}\unsure{Y|Z}\leq X,\ \unsure{\mathrm{OPS}|\mathrm{ops}}\quad
 \deleted{Z'}\unsure{Y'|Z'}\leq X'
 \end{aligned}
\]

% slanted
\margin{Axiome de \deleted{la \unsure{strate minimale}} \inserted{densité}}
\[
 \forall\,X\in M,\ \exists\,Z\in\Lambda\quad\text{tel que }\ Z\mathrel{\overset{\circ}{\ll}}X
 \tag*{$\mathrm M\Lambda\,2$}
\]

% slanted
\margin{Axiome de la strate minimale}
\[
 \begin{aligned}
 &\forall\,X\in M,\ Z\in\Lambda\quad\text{avec }\ Z\ll X,\ \exists\,X'\leq X\ \text{avec}\\
 &Z\mathrel{\overset{\circ}{\ll}}X'
 \end{aligned}
 \tag*{$\mathrm M\Lambda\,3$}
\]
\endgroup

% slanted
\margin{Pour le moment il s'agit \unsure{de|des} axiomes ne faisant pas intervenir $\ooalign{\hfil$>$\hfil\cr\hfil$<$\hfil\cr}$}
\deleted{\unsure{Il est}} Il est commode d'introduire les
\[
 \begin{aligned}
 \unsure{\mathrm{omb}|\mathrm{ends}|\mathrm{ens}}_{\Lambda}(X)
    &=\{Z\in\Lambda\mid Z\ll X\}\\
 \mathrm{omb}_{\Lambda}(X)^{\unsure{\circ|0}}
    &=\{Z\in\Lambda\mid Z\mathrel{\overset{\circ}{\ll}} X\}\,.
 \end{aligned}
\]
Les raisonnements précédents montrent la validité des deux lemmes cruciaux, que je rappelle pour mémoire :
\scanpage{41}
\mspage{40}

% slanted
\margin{cf lemme 1 p.\unsure{25|23}}
\uline{Lemme} A Soient $X,Y\in M$, supposons $\exists\,Z$ avec $X,Y\deleted{\unsure{\leq}}\ll Z$. Conditions équivalentes
\[
 \begin{aligned}
 &\text{(i)}\quad X\leq Y\\
 &\text{(ii)}\quad X\ll Y\\
 &\text{(ii bis)}\quad\unsure{\mathrm{omb}|\mathrm{ends}|\mathrm{ens}}_{\Lambda}(X)\subset\mathrm{omb}_{\Lambda}(Y)\\
 &\text{(iii)}\quad\mathrm{omb}_{\Lambda}(X)^{\circ}\subset\mathrm{omb}_{\Lambda}(Y)\\
 &\text{(iv)}\quad\mathrm{omb}_{\Lambda}(X^{\circ})\cap\mathrm{omb}_{\Lambda}(Y)\unsure{\pm|\ne}\varnothing
 \end{aligned}
\]

% slanted
\margin{cf cor 1 p.28}
\uline{Lemme} B Soit $Y\ll X$, alors $\exists\,X'$ avec $Y\mathrel{\overset{\circ}{\ll}}X'\leq X$.

% slanted
\margin{cf cor.2 p.28}
\uline{Lemme} C Soit
\diagram{\begin{tikzcd}[ampersand replacement=\&,row sep=0pt,column sep=3pt]
 Y\ll \& X\\
 \& \rotatebox{90}{$\leq$}\\
 \& X'
\end{tikzcd}}
. Conditions \unsure{équival|équiv.}
\[
 \begin{aligned}
 &\text{(i)}\quad Y\ll X'\\
 &\text{(\deleted{\unsure{ii}}\unsure{i bis|ii bis})}\quad\mathrm{omb}_{\Lambda}(Y)\subset\mathrm{omb}_{\Lambda}(X')\\
 &\text{(ii)}\quad\mathrm{omb}_{\Lambda}(Y)^{\circ}\subset\mathrm{omb}_{\Lambda}(\unsure{X'|X''})\\
 &\text{(iii)}\quad\mathrm{omb}_{\Lambda}(Y)^{\circ}\cap\mathrm{omb}_{\Lambda}(X')\ne\varnothing
 \end{aligned}
\]

% slanted
\margin{cf cor.3 p.28}
\uline{Corollaire} Mêmes notations. Cond. équivalentes
\[
 \begin{aligned}
 &\text{(i)}\quad Y\mathrel{\overset{\circ}{\ll}}X'\\
 &\text{(\deleted{\unsure{ii}}\unsure{i bis|ii bis})}\quad\mathrm{omb}_{\Lambda}(Y)^{\circ}\subset\mathrm{omb}_{\Lambda}^{\unsure{\cdot|\prime}}(X')^{\circ}\\
 &\text{(\unsure{idi|iii|ii})}\quad\mathrm{omb}_{\Lambda}(Y)^{\circ}\cap\mathrm{omb}_{\unsure{\alpha|\Lambda}}(X')^{\circ}\ne\varnothing\\
 &\text{(iv)}\quad X'\text{ est minimal \unsure{dans|de} }\widetilde X^{\deleted{\unsure{Z}}Y}
   =\{\unsure{X'|x'}\in\widetilde X\mid Z\ll X'\}\\
 &\text{(v)}\quad X'\text{ est plus petit \unsure{élément|élém} de }\widetilde X^{Y}.
 \end{aligned}
\]

Nous faisons maintenant intervenir $\unsure{\ooalign{\hfil$>$\hfil\cr\hfil$<$\hfil\cr}|\gtrless}$, et allons voir quelles propriétés de $\ooalign{\hfil$>$\hfil\cr\hfil$<$\hfil\cr}$ font marcher nos raisonnements précédents (avant de \deleted{\unsure{tout}} revenir à $\parallel$). On posera ici, pour $X,Y\in M$
\[
 X\parallel Y\ \overset{\mathrm{def}}{\Longleftrightarrow}\ X\mathrel{\ooalign{\hfil$>$\hfil\cr\hfil$<$\hfil\cr}}Y,\ \text{et }\deleted{\unsure{\nexists|\exists}}\ \widetilde X\cap\widetilde Y=\varnothing
\]

% slanted
\margin{axiome trivial pour $\ooalign{\hfil$>$\hfil\cr\hfil$<$\hfil\cr}$}
$\text{\uline{$\mathrm M\deleted{\unsure{L}}\Lambda\,0'$}}$\quad $\ooalign{\hfil$>$\hfil\cr\hfil$<$\hfil\cr}$ est réfl. et sym. Si $X\mathrel{\ooalign{\hfil$>$\hfil\cr\hfil$<$\hfil\cr}}Y$, et $X'\leq X,\ Y'\leq Y$, on a $X'\mathrel{\ooalign{\hfil$>$\hfil\cr\hfil$<$\hfil\cr}}Y'$.

% slanted
\margin{Axiome de $\Lambda$-disjonction}
$\text{\uline{$\mathrm M\Lambda\,4$}}$\quad \deleted{Soient}
\diagram{\begin{tikzcd}[ampersand replacement=\&,row sep=0pt,column sep=3pt]
 \deleted{\deleted{\unsure{X\ll}}\,\deleted{Z}\unsure{X'|Z'}\ll} \& \deleted{X}\\
 \& \deleted{\rotatebox{90}{$\leq$}}\\
 \deleted{Y'\ll} \& \deleted{\deleted{\unsure{X}}Y}
\end{tikzcd}}
\deleted{, \inserted{avec $\deleted{\unsure{Z\in\Lambda,\ Z\in|Z\in\Lambda,\ Z\ll}}\ X',Y'\in\Lambda$} Supposons $\deleted{\unsure{Z}}\unsure{X_{Y'}|X'_{Y}|Z_{Y'}|Z'_{Y}}=\varnothing$ \unsure{i.e.|tel|les} $\nexists\,\deleted{\unsure{Z}}X''\leq X'$, avec $X''\ll Y$. Alors $X'\parallel Y'$.}
